\documentclass[UTF8]{ctexart}
\usepackage[a4paper,margin=2.5cm]{geometry}
\usepackage{graphicx}
\usepackage{parskip}

\graphicspath{{./}{../chunks/}{../}}

\newcommand{\speaker}[2]{\par\medskip\noindent\textbf{#1 / #2}\par}
\newcommand{\en}[1]{\noindent\textit{#1}\par}
\newcommand{\zh}[1]{\noindent #1\par}
\newcommand{\preserveimage}[1]{%
  \begin{figure}[htbp]
    \centering
    \IfFileExists{#1}{%
      \includegraphics[width=0.9\linewidth]{#1}%
    }{%
      \fbox{\parbox{0.86\linewidth}{\centering Image reference preserved: \texttt{\detokenize{#1}}}}%
    }%
  \end{figure}%
}

\title{Republic, Book VII / 《理想国》第七卷}
\author{Plato / 柏拉图}
\date{}

\begin{document}
\maketitle

\section*{Socrates' Narration Continues / 苏格拉底继续叙述}

\speaker{SOCRATES}{苏格拉底}
\en{Next, then, compare the effect of education and that of the lack of it on our nature to an experience like this. Imagine human beings living in an underground, cavelike dwelling, with an entrance a long way up that is open to the light and as wide as the cave itself. They have been there since childhood, with their necks and legs fettered, so that they are fixed in the same place, able to see only in front of them, because their fetter prevents them from turning their heads around. Light is provided by a fire burning far above and behind them. Between the prisoners and the fire, there is an elevated road stretching. Imagine that along this road a low wall has been built--like the screen in front of people that is provided by puppeteers, and above which they show their puppets.}
\zh{那么接下来，请把教育以及缺乏教育对我们本性的影响，同下面这种经历相比较。设想有一些人住在地下洞穴般的居所里，洞口在远处上方，向光敞开，宽度与洞穴本身相当。他们自幼就在那里，颈项和腿都被镣铐束缚，因此固定在原处，只能看见面前的东西，因为束缚使他们不能转头。光来自他们身后高处燃烧的一团火。在囚徒和火之间，有一条高起的道路延伸过去。再设想沿着这条路筑起了一道矮墙--就像傀儡戏艺人置于观众面前、并在其上方展示傀儡的屏幕。}

\speaker{GLAUCON}{格劳孔}
\en{I am imagining it.}
\zh{我想象到了。}

\speaker{SOCRATES}{苏格拉底}
\en{Also imagine, then, that there are people alongside the wall carrying multifarious artifacts that project above it--statues of people and other animals, made of stone, wood, and every material. And as you would expect, some of the carriers are talking and some are silent.}
\zh{那么也请设想，墙边还有一些人，携带着各种伸出墙头的器物--用石头、木头以及各类材料制成的人像和其他动物像。正如你会预料的那样，搬运者有些在说话，有些沉默不语。}

\speaker{GLAUCON}{格劳孔}
\en{It is a strange image you are describing, and strange prisoners.}
\zh{你描绘的是一个奇特的图像，囚徒也奇特。}

\speaker{SOCRATES}{苏格拉底}
\en{They are like us. I mean, in the first place, do you think these prisoners have ever seen anything of themselves and one another besides the shadows that the fire casts on the wall of the cave in front of them?}
\zh{他们就像我们。我的意思是，首先，你认为这些囚徒除了火光投在他们面前洞壁上的自己和彼此的影子之外，还曾看见过自己或同伴的任何东西吗？}

\speaker{GLAUCON}{格劳孔}
\en{How could they, if they have to keep their heads motionless throughout life?}
\zh{如果他们一生都不得不让头保持不动，又怎么可能看见呢？}

\speaker{SOCRATES}{苏格拉底}
\en{What about the things carried along the wall? Isn't the same true where they are concerned?}
\zh{那么沿墙被搬运的那些东西呢？关于它们，情形不也是一样吗？}

\speaker{GLAUCON}{格劳孔}
\en{Of course.}
\zh{当然。}

\speaker{SOCRATES}{苏格拉底}
\en{And if they could engage in discussion with one another, don't you think they would assume that the words they used applied to the things they see passing in front of them?}
\zh{如果他们能够彼此交谈，你不认为他们会以为自己所用的词语，就是指他们眼前经过的那些东西吗？}

\speaker{GLAUCON}{格劳孔}
\en{They would have to.}
\zh{必然如此。}

\speaker{SOCRATES}{苏格拉底}
\en{What if their prison also had an echo from the wall facing them? When one of the carriers passing along the wall spoke, do you think they would believe that anything other than the shadow passing in front of them was speaking?}
\zh{如果他们的囚室还会从对面的墙壁传来回声，又会怎样？当一个沿墙走过的搬运者说话时，你认为他们会相信说话的不是眼前经过的影子，而是别的什么东西吗？}

\speaker{GLAUCON}{格劳孔}
\en{I do not, by Zeus.}
\zh{宙斯为证，我不认为会。}

\speaker{SOCRATES}{苏格拉底}
\en{All in all, then, what the prisoners would take for true reality is nothing other than the shadows of those artifacts.}
\zh{总而言之，囚徒们所当作真正实在的，不会是别的，只会是那些器物的影子。}

\speaker{GLAUCON}{格劳孔}
\en{That's entirely inevitable.}
\zh{这完全不可避免。}

\speaker{SOCRATES}{苏格拉底}
\en{Consider, then, what being released from their bonds and cured of their foolishness would naturally be like, if something like this should happen to them. When one was freed and suddenly compelled to stand up, turn his neck around, walk, and look up toward the light, he would be pained by doing all these things and be unable to see the things whose shadows he had seen before, because of the flashing lights. What do you think he would say if we told him that what he had seen before was silly nonsense, but that now--because he is a bit closer to what is, and is turned toward things that are more--he sees more correctly? And in particular, if we pointed to each of the things passing by and compelled him to answer what each of them is, don't you think he would be puzzled and believe that the things he saw earlier were more truly real than the ones he was being shown?}
\zh{那么请考虑一下：如果这样的事发生在他们身上，他们摆脱束缚、治愈愚昧，按自然情形会是什么样子。当其中一人被释放，并突然被迫站起、转动脖子、行走、抬头看向光明时，做这些事都会使他痛苦；由于光芒刺眼，他也无法看见那些过去只见其影子的东西。如果我们告诉他，他从前所见只是无谓的幻象，而现在--因为他稍微更接近存在，并转向更具有实在性的事物--他看得更正确，你认为他会怎么说？尤其是，如果我们指着每一个经过的东西，强迫他回答它是什么，你不认为他会感到困惑，并相信他早先所见比现在展示给他的东西更真实地存在吗？}

\speaker{GLAUCON}{格劳孔}
\en{Much more so.}
\zh{肯定会更这样认为。}

\speaker{SOCRATES}{苏格拉底}
\en{And if he were compelled to look at the light itself, wouldn't his eyes be pained and wouldn't he turn around and flee toward the things he is able to see, and believe that they are really clearer than the ones he is being shown?}
\zh{如果他被迫去看光本身，他的眼睛岂不会疼痛？他岂不会转身逃向那些自己能够看见的东西，并相信它们实际上比正在展示给他的那些东西更清楚吗？}

\speaker{GLAUCON}{格劳孔}
\en{He would.}
\zh{他会的。}

\speaker{SOCRATES}{苏格拉底}
\en{And if someone dragged him by force away from there, along the rough, steep, upward path, and did not let him go until he had dragged him into the light of the sun, wouldn't he be pained and angry at being treated that way? And when he came into the light, wouldn't he have his eyes filled with sunlight and be unable to see a single one of the things now said to be truly real?}
\zh{如果有人用强力把他从那里拖走，沿着崎岖、陡峭、向上的道路前行，直到把他拖进太阳之光中才放手，他岂不会因为受到这样的对待而痛苦和愤怒？当他来到光中，眼里充满阳光时，岂不会连现在所谓真正实在之物中的任何一个都看不见吗？}

\speaker{GLAUCON}{格劳孔}
\en{No, he would not be able to--at least not right away.}
\zh{是的，他不能，至少不会立刻能。}

\speaker{SOCRATES}{苏格拉底}
\en{He would need time to get adjusted, I suppose, if he is going to see the things in the world above. At first, he would see shadows most easily, then images of men and other things in water, then the things themselves. From these, it would be easier for him to go on to look at the things in the sky and the sky itself at night, gazing at the light of the stars and the moon, than during the day, gazing at the sun and the light of the sun.}
\zh{我想，如果他要看见上方世界中的事物，就需要时间来适应。起初，他最容易看见影子；接着，看见水中人和其他事物的映像；然后，才看见事物本身。由此再进一步，他在夜间观看天上的事物和天空本身，凝视星月之光，会比在白昼凝视太阳及太阳之光更容易。}

\speaker{GLAUCON}{格劳孔}
\en{Of course.}
\zh{当然。}

\speaker{SOCRATES}{苏格拉底}
\en{Finally, I suppose, he would be able to see the sun--not reflections of it in water or some alien place, but the sun just by itself in its own place--and be able to look at it and see what it is like.}
\zh{最后，我想，他就能够看见太阳--不是它在水中或异处的倒影，而是太阳在自身所在之处、就其自身而言的样子--并且能够直视它，看清它是什么样。}

\speaker{GLAUCON}{格劳孔}
\en{Necessarily.}
\zh{必然如此。}

\speaker{SOCRATES}{苏格拉底}
\en{After that, he would already be able to conclude about it that it provides the seasons and the years, governs everything in the visible world, and is in some way the cause of all the things that he and his fellows used to see.}
\zh{在那之后，他便已经能够关于太阳得出结论：它提供季节和岁月，统摄可见世界中的一切，并且在某种意义上，是他和同伴们过去所见一切的原因。}

\speaker{GLAUCON}{格劳孔}
\en{That would clearly be his next step.}
\zh{显然，这会是他的下一步。}

\speaker{SOCRATES}{苏格拉底}
\en{What about when he reminds himself of his first dwelling place, what passed for wisdom there, and his fellow prisoners? Don't you think he would count himself happy for the change and pity the others?}
\zh{那么，当他回想自己最初的住处、那里所谓的智慧以及那些同伴囚徒时，又会怎样？你不认为他会为自己的改变而庆幸，并怜悯其他人吗？}

\speaker{GLAUCON}{格劳孔}
\en{Certainly.}
\zh{当然会。}

\speaker{SOCRATES}{苏格拉底}
\en{And if there had been honors, praises, or prizes among them for the one who was sharpest at identifying the shadows as they passed by; and was best able to remember which usually came earlier, which later, and which simultaneously; and who was thus best able to prophesy the future, do you think that our man would desire these rewards or envy those among the prisoners who were honored and held power? Or do you think he would feel with Homer that he would much prefer ``to work the earth as a serf for another man, a man without possessions of his own,'' and go through any sufferings, rather than share their beliefs and live as they do?}
\zh{如果他们中间曾有荣誉、赞扬或奖赏，颁给那个最敏锐地辨认经过影子的人；颁给那个最能记住哪些影子通常在先、哪些在后、哪些同时出现的人；也就是颁给那个最能预言未来的人，那么你认为我们这个人会渴望这些奖赏，或嫉妒那些在囚徒中受尊荣、掌权的人吗？还是说，你认为他会同荷马一起觉得，自己宁愿“替另一个人，一个没有自己财产的人，像雇工那样耕作土地”，并忍受任何痛苦，也不愿分享他们的信念、过他们那样的生活？}

\speaker{GLAUCON}{格劳孔}
\en{Yes, I think he would rather suffer anything than live like that.}
\zh{是的，我认为他宁愿忍受任何事，也不愿那样生活。}

\speaker{SOCRATES}{苏格拉底}
\en{Consider this too, then. If this man went back down into the cave and sat down in his same seat, wouldn't his eyes be filled with darkness, coming suddenly out of the sun like that?}
\zh{那么也请考虑这一点。如果这个人再下到洞穴中，坐回原来的座位，他这样突然从太阳下回来，眼睛岂不会充满黑暗吗？}

\speaker{GLAUCON}{格劳孔}
\en{Certainly.}
\zh{当然会。}

\speaker{SOCRATES}{苏格拉底}
\en{Now, if he had to compete once again with the perpetual prisoners in recognizing the shadows, while his sight was still dim and before his eyes had recovered, and if the time required for readjustment was not short, wouldn't he provoke ridicule? Wouldn't it be said of him that he had returned from his upward journey with his eyes ruined, and that it is not worthwhile even to try to travel upward? And as for anyone who tried to free the prisoners and lead them upward, if they could somehow get their hands on him, wouldn't they kill him?}
\zh{现在，如果他在视力仍然昏暗、眼睛还没有恢复的时候，又不得不再次和那些终身囚徒比赛辨认影子，而且重新适应所需的时间并不短，他岂不会招来嘲笑？人们岂不会说，他完成上行之旅回来，眼睛已经坏了，甚至连尝试上行都不值得？至于任何试图释放囚徒并把他们带向上方的人，如果他们能设法抓住他，岂不会杀了他吗？}

\speaker{GLAUCON}{格劳孔}
\en{They certainly would.}
\zh{他们一定会。}

\speaker{SOCRATES}{苏格拉底}
\en{This image, my dear Glaucon, must be fitted together as a whole with what we said before. The realm revealed through sight should be likened to the prison dwelling, and the light of the fire inside it to the sun's power. And if you think of the upward journey and the seeing of things above as the upward journey of the soul to the intelligible realm, you won't mistake my intention--since it is what you wanted to hear about. Only the god knows whether it is true. But this is how these phenomena seem to me: in the knowable realm, the last thing to be seen is the form of the good, and it is seen only with toil and trouble. Once one has seen it, however, one must infer that it is the cause of all that is correct and beautiful in anything, that in the visible realm it produces both light and its source, and that in the intelligible realm it controls and provides truth and understanding; and that anyone who is to act sensibly in private or public must see it.}
\zh{亲爱的格劳孔，这个图像必须作为一个整体，同我们先前所说的东西配合起来理解。由视觉显现的领域，应当比作囚居之所；其中火的光，应当比作太阳的力量。而如果你把上行之旅以及观看上方事物，理解为灵魂上升到可知领域的旅程，你就不会误解我的意思--因为这正是你想听的。至于它是否为真，只有神知道。但在我看来，情形是这样：在可知领域中，最后被看见的是善之理念（善的理型），而且只有经过辛劳困苦才能看见。然而，一旦看见它，就必须推知：它是一切事物中正当者与美好者的原因；在可见领域中，它产生光及光源；在可知领域中，它主宰并提供真理和理解；并且，任何想要在私人事务或公共事务中明智行事的人，都必须看见它。}

\speaker{GLAUCON}{格劳孔}
\en{I agree, so far as I am able.}
\zh{就我所能理解的而言，我同意。}

\speaker{SOCRATES}{苏格拉底}
\en{Come on, then, and join me in this further thought: you should not be surprised that the ones who get to this point are not willing to occupy themselves with human affairs, but that, on the contrary, their souls are always eager to spend their time above. I mean, that is surely what we would expect, if indeed the image I described before is also accurate here.}
\zh{那么，请继续同我一起思考这一点：那些达到这一境地的人不愿忙于人间事务，反而其灵魂总是渴望在上方停留，对此你不应感到惊讶。我的意思是，如果我先前描述的图像在这里也确实恰当，这当然正是我们会预期的情形。}

\speaker{GLAUCON}{格劳孔}
\en{It is what we would expect.}
\zh{这正是我们会预期的。}

\speaker{SOCRATES}{苏格拉底}
\en{What about when someone, coming from looking at divine things, looks to the evils of human life? Do you think it is surprising that he behaves awkwardly and appears completely ridiculous, if--while his sight is still dim and he has not yet become accustomed to the darkness around him--he is compelled, either in the courts or elsewhere, to compete about the shadows of justice, or about the statues of which they are the shadows; and to dispute the way these things are understood by people who have never seen justice itself?}
\zh{那么，当一个人刚观看过神圣事物，转而察看人类生活的诸种恶时，又会怎样？如果--在他的视力仍然昏暗、还没有习惯周围黑暗的时候--他被迫在法庭或别处，就正义的影子，或就投下那些影子的雕像展开争竞，并且同那些从未见过正义本身的人争论这些东西应如何理解，你认为他举止笨拙、显得十分可笑，会令人惊讶吗？}

\speaker{GLAUCON}{格劳孔}
\en{It is not surprising at all.}
\zh{这根本不令人惊讶。}

\begin{center}
\textit{[\dots]}
\end{center}

\preserveimage{images/c71e19475dd43f2902a91e328edeb2062a34b8a6f56720da0301812d3489866e.jpg}
\preserveimage{images/aaa4581cac55b98fb3bd196278730238d2964fc259e31265903b352a51812148.jpg}
\preserveimage{images/2509989365e4ca64c8315dcf70650facf7be0779ef7a983d4ed5c3e65ad553a6.jpg}

\end{document}
